\section{问题分析}

\subsection{问题一的分析}
问题一本质上是一个 XXXX 问题。考虑到 XXXX，我们计划采用 XXXX 方法进行求解，主要思路为：先 XXXX，再 XXXX，最后 XXXX。

\subsection{问题二的分析}
问题二在问题一基础上增加了 XXXX 约束 / 因素，需要 XXXX。

\subsection{问题三的分析}
问题三 XXXX。

\subsection{整体求解思路}
本文整体求解思路如图~\ref{fig:framework} 所示。

\begin{figure}[H]
    \centering
    \begin{tikzpicture}[
        node distance=1.2cm,
        box/.style={rectangle, rounded corners, draw=black, thick,
                    minimum width=6em, minimum height=2em, align=center, font=\small},
        arrow/.style={-Stealth, thick}
    ]
        \node[box, fill=blue!10]   (p)  {题目分析};
        \node[box, fill=blue!10,  below=of p]  (a)  {模型假设};
        \node[box, fill=green!15, below=of a]  (m1) {问题一建模};
        \node[box, fill=green!15, below=of m1] (m2) {问题二建模};
        \node[box, fill=green!15, below=of m2] (m3) {问题三建模};
        \node[box, fill=orange!20, below=of m3] (v)  {模型检验与推广};
        \node[box, fill=red!10,    below=of v]  (c)  {结论};

        \draw[arrow] (p)  -- (a);
        \draw[arrow] (a)  -- (m1);
        \draw[arrow] (m1) -- (m2);
        \draw[arrow] (m2) -- (m3);
        \draw[arrow] (m3) -- (v);
        \draw[arrow] (v)  -- (c);
    \end{tikzpicture}
    \caption{本文整体求解框架}
    \label{fig:framework}
\end{figure}
